\documentclass[10pt,a4paper]{article}
\usepackage[margin=24mm]{geometry}
\usepackage{amsmath,amssymb,amsthm,booktabs,longtable,hyperref}
\setlength{\emergencystretch}{2em}
\newtheorem{theorem}{Theorem}\newtheorem{lemma}[theorem]{Lemma}
\newtheorem{proposition}[theorem]{Proposition}
\newcommand{\Aut}{\operatorname{Aut}}
\newcommand{\Sp}{\operatorname{Sp}}\newcommand{\SL}{\operatorname{SL}}
\newcommand{\SO}{\operatorname{SO}}\newcommand{\GL}{\operatorname{GL}}
\newcommand{\PSL}{\operatorname{PSL}}\newcommand{\PGL}{\operatorname{PGL}}
\newcommand{\Inn}{\operatorname{Inndiag}}
\newcommand{\nuu}{\nu_2}
\title{Involutions outside the reversal and subgroup obstructions}
\author{Human Compression project}\date{}
\begin{document}\maketitle

\begin{theorem}\label{main}
Let $G$ be a finite simple exceptional group of Lie type, other than
$G_2(2)'\cong\operatorname{PSU}_3(3)$ and
${}^2G_2(3)'\cong\PSL_2(8)$. There are $x,t\in G$ such that
\[
 |t|=2,\qquad \langle x,t\rangle=G,
 \qquad \nexists\alpha\in\Aut(G):\ (x^\alpha,t^\alpha)=(x^{-1},t).
\]
\end{theorem}

This is the exceptional-group part of the positive existence theorem
of Leemans and Liebeck~\cite{LL}. The already reviewed finite certificates
cover $G_2(3)$, $G_2(4)$ and the Tits group ${}^2F_4(2)'$.
The argument below treats every remaining parameter.
The passage from these generators to an abstract chiral polyhedron is
the retained Schulte--Weiss reconstruction theorem. In a nonabelian simple
group the cyclic-subgroup intersection required there is automatic:
$\langle x\rangle\cap\langle x^{-1}t\rangle$ centralizes both generators,
and is therefore trivial.

\section{A counting criterion with the outer automorphisms included}

For a finite group $H$, write $i_2(H)$ for its number of elements of
order two and $\nuu(H)=i_2(H)+1$. Thus $\nuu$ counts solutions of $h^2=1$,
including the identity. If $K\lhd H$, then
\begin{equation}\label{extension}
 \nuu(H)\le |K|\nuu(H/K).
\end{equation}
Indeed, each solution in $H$ maps to a solution in the quotient, and a
fibre contains at most $|K|$ elements. If $K$ is central of odd order, each
such fibre contains at most one solution: any two differ by a central
element whose square is one. We will use both forms, keeping the identity
in every estimate involving an extension.

\begin{lemma}\label{cosets}
Let $G$ be centreless, let $G\le J\lhd A=\Aut(G)$, and suppose
$|x|>2$ and $C_J(x)=T$. Put
\[
 \kappa=|\{aJ\in A/J:(aJ)^2=J\}|.
\]
Suppose $\mu$ bounds $i_2(C_G(\alpha))$ for every involution $\alpha\in A$
which sends $x$ to $x^{-1}$. Suppose every proper subgroup containing $x$
lies in either $N_G(\langle x\rangle)$ or one of the listed maximal
subgroups $M_1,\ldots,M_s$. If
\begin{equation}\label{criterion}
 i_2(G)>\kappa|T|\mu+
       |N_G(\langle x\rangle)|+\sum_{j=1}^s i_2(M_j),
\end{equation}
then $x$ has an involutory partner as in Theorem~\ref{main}.
\end{lemma}
\begin{proof}
Fix a coset of $J$ containing a reverser $a$. Every other reverser in
that coset differs from $a$ by an element of $C_J(x)$; this follows by
applying their quotient to $x$. The set of reversers in that coset thus
has at most $|T|$ elements. An involutory reverser has image of square one
in $A/J$, so there are at most $\kappa|T|$ of them in total.
The first term of~\eqref{criterion} consequently bounds the involutions
fixed by any involutory reverser. The remaining terms bound all
involutions which fail to generate $G$ with $x$. Choose $t$ outside this
union. If an automorphism reverses $x$ and fixes $t$, its square fixes
the generating pair and is the identity. It is not the identity, because
$|x|>2$. It would therefore be one of the excluded involutory reversers,
a contradiction.
\end{proof}

The source uses the corresponding single-coset estimate. Our form only
requires the centralizer inside $J$; field and graph automorphisms are
counted explicitly. It also charges the torus normalizer explicitly.

\section{The tori and the retained subgroup classification}

Let $\mathbf G$ be the adjoint algebraic group, let $F$ be its defining
endomorphism, and put $J=\mathbf G^F$, $G=J'$. Choose the cyclic torus
$T_0=\langle x\rangle$ in Table~\ref{tori}. Its existence and cyclicity
are the torus results collected in Kantor--Seress~\cite[Section~2]{KS}.
The index of $T_0$ in the corresponding torus $T$ of $J$ is
$d=(2,q-1)$ for $E_7$, $d=(3,q-\epsilon)$ for $E_6^\epsilon$, and $d=1$
otherwise. In the table $D$ is the exact order of $T$.

\begin{table}[ht]\centering\small
\begin{tabular}{lll}
\toprule
$G$&$D$&other maximal overgroups of $T_0$\\\midrule
$E_8(q)$&$\Phi_{30}(q)$&none\\
$E_7(q),\ q\ge3$&$\Phi_{18}(q)(q+1)$&one ${}^2E_6(q).(q+1)/d$\\
$E_7(2)$&$129$&one $\operatorname{PSU}_8(2).2$\\
$E_6(q)$&$\Phi_9(q)$&one $\PSL_3(q^3).3$\\
$ {}^2E_6(q)$&$\Phi_{18}(q)$&one $\operatorname{PSU}_3(q^3).3$\\
$F_4(q),\ q\ge3$&$\Phi_{12}(q)$&at most two ${}^3D_4(q).3$\\
$F_4(2)$&$17$&two $\Sp_8(2)$\\
$G_2(q),\ q\ge5$&$\Phi_6(q)$&one $\operatorname{SU}_3(q).2$\\
$ {}^3D_4(q)$&$\Phi_{12}(q)$&none\\
$ {}^2F_4(q),\ q\ge8$&$q^2+\sqrt{2q^3}+q+\sqrt{2q}+1$&none\\
$ {}^2G_2(q),\ q\ge27$&$q+\sqrt{3q}+1$&none\\
$ {}^2B_2(q),\ q\ge8$&$q+\sqrt{2q}+1$&none\\\bottomrule
\end{tabular}
\caption{The chosen tori and their other maximal overgroups.}
\label{tori}
\end{table}

The overgroup information, including the numbers of overgroups containing
the chosen torus rather than just their conjugacy types, is retained from
Guralnick--Kantor~\cite[Propositions~6.1--6.2 and their proofs]{GK}.
Their small-case Table~IV supplies $F_4(2)$ and $E_7(2)$.
Central quotients in Table~\ref{tori} are understood in their natural
embeddings; the estimates below allow a covering kernel wherever needed.
We are not reproving the classification of maximal subgroups.

We do prove the centralizer interface required for Lemma~\ref{cosets}.
This avoids applying a general torus-centralizer statement outside its
stated field-size range.

\begin{lemma}\label{regular}
For the choices in Table~\ref{tori},
$C_{\mathbf G}(x)=\mathbf T$ is the connected algebraic torus containing
$T$, and $C_J(x)=T$. Moreover,
\begin{equation}\label{normalizer}
 |N_G(\langle x\rangle)|\le |W|D,
\end{equation}
where $W$ is the absolute Weyl group.
\end{lemma}
\begin{proof}
Let $Q$ be the character lattice of the adjoint torus, which is the root
lattice. For an ordinary twisted or untwisted group the endomorphism on
this lattice is $F^*=qu$, with $u$ an isometry. A root $\alpha$ trivial on
$T_0$ gives a character of $T/T_0$, and hence
\[
 d\alpha=(qu-I)v\quad\hbox{for some }v\in Q.
\]
Consequently
\begin{equation}\label{norm}
 \|v\|\le\frac{d\|\alpha\|}{q-1}.
\end{equation}
Equality requires $uv=v$. In the simply laced cases all roots have the
length of a shortest nonzero lattice vector. The chosen torus types have
no eigenvalue one: their characteristic polynomials are
$\Phi_{30}$, $\Phi_{18}\Phi_2$ (or $\Phi_{14}\Phi_2$ for $E_7(2)$),
$\Phi_9$, $\Phi_{18}$ and $\Phi_{12}$, respectively.
For $E_8,E_7,E_6^\epsilon,{}^3D_4$ the right side
of~\eqref{norm} is at most this minimum length, except for ${}^2E_6(2)$.
If equality could occur, the absence of a fixed vector makes the actual
inequality strict. Thus no root is trivial on $T_0$.
For $F_4(q)$ with $q\ge3$ and $G_2(q)$ with $q\ge5$, the ratios of longest
to shortest root length are $\sqrt2$ and $\sqrt3$, so~\eqref{norm}
is again strictly below the minimum lattice length.

For the Suzuki and Ree endomorphisms, $F^*/\sqrt q$ is an isometry.
Replacing $q-1$ by $\sqrt q-1$ proves the same assertion:
at the smallest allowed parameters the relevant upper ratios are
$\sqrt2/(\sqrt8-1)<1$ and $\sqrt3/(\sqrt{27}-1)<1$.
This includes ${}^2F_4(8)$.

The two remaining lattice checks are finite. We use simple-root bases
with the $E_6$ diagram edges
$(1,2),(2,3),(3,4),(4,5),(3,6)$ and the $F_4$ Cartan matrix
\[
 \begin{pmatrix}2&-1&0&0\\-1&2&-2&0\\0&-1&2&-1\\0&0&-1&2\end{pmatrix}.
\]
The supplied certificate constructs all $72$ and $48$ roots by their
simple reflections. It supplies a Weyl matrix $w$ of characteristic
polynomial $\Phi_9$ in $E_6$ and one of polynomial $\Phi_8$ in $F_4$.
For every root, respectively,
\[
 (-2w-I)^{-1}(3\alpha)\notin\mathbb Z^6,
 \qquad (2w-I)^{-1}\alpha\notin\mathbb Z^4.
\]
The determinants have absolute values $57$ and $17$. These are the
required two torus types, and the displayed tests prove regularity in
${}^2E_6(2)$ and $F_4(2)$. Thus in all cases
$C_{\mathbf G}(x)^\circ=\mathbf T$.

We check connectedness as well. For a central isogeny
$\widetilde{\mathbf G}\to\mathbf G$ from the simply connected group,
the component group of a semisimple centralizer embeds in the central
kernel: an element centralizing $x$ takes a lift $\tilde x$ to
$z\tilde x$, and the kernel of this map is connected. If $|x|$ is coprime
to the order of the kernel, then $z^{|x|}=1$ forces $z=1$.
This applies to $E_6^\epsilon$ because $\Phi_9(q)/d$ and
$\Phi_{18}(q)/d$ are prime to three; when three divides the numerator
its valuation is exactly one. It applies to ${}^3D_4$ because its torus
has odd order and the covering kernel is a two-group. The other cases
have trivial central kernel on geometric points, except for odd-characteristic
$E_7$.

For $E_7$, the exact Weyl calculation gives
$C_W(u)=\langle u\rangle$ of order $18$, or $14$ at $q=2$.
In the rational normalizer $u$ acts on $T$ by the inverse of the
$q$-power map. For $q\ge3$, $\Phi_{18}(q)$ divides $|T_0|$, and $q$ has
order $18$ modulo $\Phi_{18}(q)$: its order divides $18$,
$q^9\equiv-1$, and $q^6\equiv q^3-1\not\equiv1$ modulo that number.
These exclude every proper divisor of $18$. At $q=2$, the order of two
modulo $129$ is $14$. Thus no nonidentity Weyl-normalizer element fixes
$x$. If its algebraic centralizer were disconnected, its component group
would have order two. Frobenius acts trivially on that group, and Lang's
theorem for the connected torus would give a rational representative of
its nontrivial component. That contradicts the faithful normalizer action.
The centralizer is therefore connected here too. The certificate checks
both Weyl centralizers, not merely the orders of chosen elements.

Finally, a normalizer of $\langle x\rangle$ normalizes its algebraic
centralizer $\mathbf T$. Its quotient by its intersection with $\mathbf T$
embeds in $W$, while that intersection has at most $D$ elements.
This proves~\eqref{normalizer}.
\end{proof}

\section{Involutions which reverse the torus generator}

Write $q=p^f$. The standard structure of outer automorphisms gives
the following bounds on the square-one elements of $\Aut(G)/J$:
\[
\begin{array}{c|rrrrrrrrr}
G&E_8&E_7&E_6^\epsilon&F_4&G_2&{}^3D_4&{}^2F_4&{}^2G_2&{}^2B_2\\\hline
\kappa&2&2&4&2&2&2&1&1&1.
\end{array}
\]
The quotient is cyclic apart from the untwisted $E_6$ graph factor,
where it is $C_f\times C_2$. In the exceptional graph-field cases
$F_4$ in characteristic two and $G_2$ in characteristic three it is
cyclic of order $2f$. For ${}^3D_4$ it is cyclic of order $3f$, so
$\kappa=1$ at $q=2,3$. In the three Suzuki--Ree families the field
degree is odd and there is no nontrivial element of square one.

Suppose first that an inner-diagonal involution reverses $x$.
For every type other than $E_6$, the longest Weyl element is $-1$.
An algebraic representative $n$ of it can be chosen to have order two;
all elements in $\mathbf Tn$ then have order two and are conjugate
over the algebraic closure, since squaring on the torus is surjective.
There is a rational representative by Lang's theorem for $\mathbf T$.
If $a$ is any other rational reverser, $an^{-1}$ centralizes $x$ and hence
lies in $T$, by Lemma~\ref{regular}. Thus all such reversers belong to
this algebraic involution class.

For $E_6$ the Weyl centralizer of the $\Phi_9$ torus element has order
nine, as checked in the root certificate. The twisted torus uses its
negative; its Weyl centralizer is the same group. Since the torus order
is odd, $N_J(T)$ has odd order. There is consequently no inner-diagonal
involutory reverser in this case.

The involution centralizer classification is a retained input here.
For odd characteristic it is summarized in~\cite[Section~2, Table~5]{BT};
for even characteristic we use Aschbacher--Seitz~\cite[Sections~13--18]{AS}
with their corrections~\cite{ASerr}. Applied to the lift of $-1$, it gives
the centralizer types in Table~\ref{centralizers}. For a covering group
or an adjoint group the indicated small central kernels are allowed.
The powers in brackets denote the order of the normal unipotent subgroup.

\begin{table}[ht]\centering\small
\begin{tabular}{lll}
\toprule
type&odd characteristic&even characteristic\\\midrule
$E_8$&$D_8$&$[q^{84}].\Sp_8(q)$\\
$E_7$&$A_7^\epsilon.2$&$[q^{42}].\Sp_6(q)$\\
$F_4$&$(A_1C_3).2$&$[q^{18}].\SL_2(q)^2$\\
$G_2$&$(A_1A_1).2$&$[q^3].\SL_2(q)$\\
$ {}^3D_4$&$(A_1(q)A_1(q^3)).2$&$[q^9].\SL_2(q)$\\
$ {}^2F_4$&---&$[q^9].\SL_2(q)$\\
$ {}^2G_2$&$\PSL_2(q)\times2$&---\\
$ {}^2B_2$&---&a Sylow two-subgroup\\\bottomrule
\end{tabular}
\caption{Retained centralizer types for a lift of the inverting Weyl element.}
\label{centralizers}
\end{table}

To estimate their involutions we retain the uniform theorem of
Lawther--Liebeck--Seitz~\cite[Proposition~1.3]{LLS}:
\begin{equation}\label{LLS}
 i_2(\Aut(H))<2(q^{M_H}+q^{M_H-1}).
\end{equation}
Here $M_H=N+r$ for the absolute root system, and half that value for a
Suzuki--Ree group. Since the upper bound is an integer, the same expression
is an upper bound for $\nuu(\Aut(H))$. This observation is needed in
using~\eqref{extension}.

For $E_8$ in odd characteristic a central kernel of order at most four
reduces the centralizer to a subgroup of $\Aut(P\Omega_{16}^+(q))$.
Thus $\nuu\le8(q^{64}+q^{63})$, which is below
$2q^{84}(q^{20}+q^{19})$. For odd $E_7$, a central kernel of order at most
eight gives $\nuu\le16(q^{35}+q^{34})$, below
$2q^{42}(q^{12}+q^{11})$. In even characteristic~\eqref{extension}
gives the latter two expressions directly.

For odd $q$, $\nuu(\PGL_2(q))=q^2+1$. The central products in the next
three rows have kernel at most two, and their quotients embed in the
products of the corresponding inner-diagonal factors. They give
\[
\begin{array}{c|c}
F_4&4(q^2+1)(q^{12}+q^{11})<q^{22}\\
G_2&2(q^2+1)^2<q^5\\
{}^3D_4&2(q^2+1)(q^6+1)<q^{11}.
\end{array}
\]
These inequalities hold for odd $q\ge3$. In even characteristic
$\nuu(\SL_2(q))=q^2$; multiplication by the indicated unipotent orders
gives $q^{22},q^5,q^{11}$, and also $q^{11}$ for ${}^2F_4$.
For the small Ree group,
$i_2(\PSL_2(q)\times2)=q(q-1)+1<2q^2$.
For the Suzuki group the acquired Sylow-subgroup calculation gives
$i_2(C_G(n))=q-1$.

The outer involutions must still be included. Proposition~1.1 of~\cite{LLS}
gives their fixed groups. A field involution has a group of the same
type over $s$ with $s^2=q$, allowing inner-diagonal extensions.
Formula~\eqref{LLS} bounds its involutions by $4q^{M_G/2}$.
This fits the preceding bounds for $E_8,E_7,F_4,G_2,{}^3D_4$.
For $E_6^\epsilon$ a field or graph-field involution is bounded by
$4q^{21}$. A graph involution has fixed group $F_4(q)$, $C_4(q)$,
or an involution centralizer in $F_4(q)$; in every case
$2(q^{28}+q^{27})$ is sufficient.
The exceptional graph-field involutions of $F_4$ and $G_2$ have fixed
groups ${}^2F_4(q)$ and ${}^2G_2(q)$ respectively. Their bounds
$2(q^{14}+q^{13})$ and $2(q^4+q^3)$ fit $q^{22}$ and $q^5$.
At $F_4(2)$ the fixed group is ${}^2F_4(2)$, which embeds by conjugation
in the automorphism group of its simple derived Tits group;
Proposition~1.3 of~\cite{LLS} still gives the bound $49152<2^{22}$.
No field involution occurs in the three Suzuki--Ree families themselves.

We have therefore proved that the quantities $C$ in Table~\ref{bounds}
bound $i_2(C_G(\alpha))$ for \emph{every} involutory reverser, not only
the inner-diagonal ones.

\section{Lower bounds from one explicit conjugacy class}

We also make the constants in the lower bounds explicit. Let $J$ again
denote the inner-diagonal group. Proposition~2.3, Remark~2.2 and Table~5
of~\cite{BT} provide the involutions in Table~\ref{classes}. Each lies in
$G$, and its $J$-class is its $G$-class. Thus $|J|/|C_J(t)|$ is a lower
bound for $i_2(G)$. These are retained class orders, not estimates inferred
from asymptotic notation.

\begin{table}[ht]\centering\small
\begin{tabular}{lll}
\toprule
type&$|C_J(t)|$, odd characteristic&$|C_J(t)|$, even characteristic\\\midrule
$E_8$&$|\SO_{16}^+(q)|$&$q^{84}|\Sp_8(q)|$\\
$E_7$&$2|\SL_8^\epsilon(q)|$, $q\equiv\epsilon\pmod4$
 &$q^{42}|\Sp_6(q)|$\\
$E_6^\epsilon$&$|\SL_2(q)|\,|\SL_6^\epsilon(q)|$
 &$q^{27}|\SL_2(q)|\,|\SL_3^\epsilon(q)|$\\
$F_4$&$|\SL_2(q)|\,|\Sp_6(q)|$&$q^{18}|\SL_2(q)|^2$\\
$G_2$&$|\SL_2(q)|^2$&$q^3|\SL_2(q)|$\\
$ {}^3D_4$&$|\SL_2(q)|\,|\SL_2(q^3)|$&$q^9|\SL_2(q)|$\\
$ {}^2F_4$&---&$q^9|\SL_2(q)|$\\
$ {}^2G_2$&$|\SL_2(q)|$&---\\
$ {}^2B_2$&---&$q^2$\\\bottomrule
\end{tabular}
\caption{Centralizer orders of the selected lower-bound classes.}
\label{classes}
\end{table}

Substitution of the standard group-order products gives
\begin{equation}\label{lower}
 i_2(G)>q^a/c,
 \quad
\begin{array}{c|rrrrrrrrr}
G&E_8&E_7&E_6^\epsilon&F_4&G_2&{}^3D_4&{}^2F_4&{}^2G_2&{}^2B_2\\\hline
a&128&70&40&28&8&16&14&4&3\\
c&2&4&2&2&2&2&2&2&2.
\end{array}
\end{equation}
Here is an exact way to check every constant. Form the integer polynomial
$P(q)=c|J|-q^a|C_J(t)|$ from the row in Table~\ref{classes}. Expand
$P(z+q_0)$, where $q_0$ is the smallest allowed parameter of that parity.
In all eighteen cases every coefficient is nonnegative and the constant
coefficient is positive. The certificate stores every coefficient, using
integer arithmetic; this proves positivity for all real $q\ge q_0$ and
hence for every required prime power. The order products used are
\[
 |\SL_n^\epsilon(q)|=q^{n(n-1)/2}\prod_{i=2}^n(q^i-\epsilon^i),
 \quad |\Sp_{2m}(q)|=q^{m^2}\prod_{i=1}^m(q^{2i}-1),
\]
and the analogous displayed exceptional order products in the certificate.
The correction to the ${}^2E_6$ centralizer in~\cite{ASerr} is respected.
For the Suzuki groups we will use the stronger exact value
$(q-1)(q^2+1)$, obtained from their single involution class.

\section{The final inequalities}

The column $B$ in Table~\ref{bounds} bounds the sum of the involution
counts of the other maximal overgroups in Table~\ref{tori}.
For $E_7$, a central fibre of size at most $q+1$ over
$\Aut({}^2E_6(q))$ gives the stated expression. The special $E_7(2)$
overgroup is bounded by $2(2^{35}+2^{34})$, which is smaller.
For $E_6^\epsilon$, the extension of index three contributes no new
involutions, and any central covering kernel is odd; apply the strengthened
form of~\eqref{extension} and~\eqref{LLS} over $q^3$.
For $F_4$ there are at most two overgroups, each in
$\Aut({}^3D_4(q))$. For $G_2$, the possible central kernel of
$\operatorname{SU}_3(q).2$ has order at most three, giving
$6(q^5+q^4)$. These estimates deliberately allow every indicated kernel.

\begin{table}[ht]\centering\small
\begin{tabular}{llll}
\toprule
$G$&$C$&$B$&$|W|$\\\midrule
$E_8$&$2q^{84}(q^{20}+q^{19})$&$0$&$696729600$\\
$E_7$&$2q^{42}(q^{12}+q^{11})$&$2(q+1)(q^{42}+q^{41})$&$2903040$\\
$E_6^\epsilon$&$2(q^{28}+q^{27})$&$2(q^{15}+q^{12})$&$51840$\\
$F_4,\ q\ge3$&$q^{22}$&$4(q^{16}+q^{15})$&$1152$\\
$G_2,\ q\ge5$&$q^5$&$6(q^5+q^4)$&$12$\\
$ {}^3D_4$&$q^{11}$&$0$&$192$\\
$ {}^2F_4$&$q^{11}$&$0$&$1152$\\
$ {}^2G_2$&$2q^2$&$0$&$12$\\
$ {}^2B_2$&$q-1$&$0$&$8$\\\bottomrule
\end{tabular}
\caption{Bounds charged to the two kinds of obstruction.}
\label{bounds}
\end{table}

For $E_6^\epsilon$ we may replace $D$ by the larger $\Phi_9(q)$,
and for $E_7(2)$ by $\Phi_{18}(2)\Phi_2(2)=171$.
Let $L$ be the right side of~\eqref{lower}, using the exact Suzuki value
instead. We verify
\begin{equation}\label{finalineq}
 L-\kappa DC-B-|W|D>0.
\end{equation}
For the ordinary groups expand the left side at $q=q_0+z$, with
$q_0=2$ for $E_8,E_7,E_6^\epsilon$, $q_0=3$ for $F_4$, $q_0=5$ for
$G_2$, and $q_0=4$ for ${}^3D_4$.
For ${}^2F_4$ and ${}^2B_2$ put $q=2s^2$ and expand at $s=2+z$;
for ${}^2G_2$ put $q=3s^2$ and expand at $s=3+z$.
Every coefficient is nonnegative and the constant is positive in all
nine cases. The supplied rational-arithmetic calculation records the
full polynomials and their translated coefficients, rather than sampling
field sizes.

At ${}^3D_4(2)$ and ${}^3D_4(3)$ use $\kappa=1$.
The exact margins in~\eqref{finalineq} are $3648$ and $17155227/2$.
At $F_4(2)$ the selected class has size $350859600$; the two overgroups
give $B\le4(2^{20}+2^{19})$. With $D=17$, $C=2^{22}$ and $\kappa=2$
the margin is $201942224$.
These calculations cover every remaining small parameter.
Lemma~\ref{cosets}, together with~\eqref{normalizer}, now proves
Theorem~\ref{main}.

\section{Attribution and verification scope}

The existence result belongs to~\cite{LL}. The choice of tori and the
strategy of excluding maximal subgroups and involution centralizers are
also borrowed from that proof. The subgroup classification, conjugacy
class orders, outer-automorphism classification and uniform upper bound
remain explicit external premises, cited where used. They are not new
theorems of this project.

The additions here are the outer-coset criterion, the root-lattice
centralizer argument including the finite checks, the explicit normalizer
charge, and the complete constant inequalities. The latter distinguish
$i_2$ from $\nu_2$ before passing through an extension. No asymptotic
estimate supplies an unverified finite constant. The arithmetic certificates
prove precisely the polynomial and lattice assertions they state;
they do not prove the retained classification theorems.
The full L045 admission additionally requires its classical, alternating
and sporadic components. No new conjecture resolution, Lean verification
or compression measurement is claimed by this component.

\begin{thebibliography}{9}\small
\bibitem{LL} D. Leemans and M. W. Liebeck, \emph{Chiral polyhedra and
finite simple groups}, arXiv:1603.07713v3, Sections 2--3.
\bibitem{GK} R. M. Guralnick and W. M. Kantor, \emph{Probabilistic
generation of finite simple groups}, J. Algebra 234 (2000), 743--792.
\bibitem{KS} W. M. Kantor and \`A. Seress, \emph{Prime power graphs
for groups of Lie type}, J. Algebra 247 (2002), 370--434.
\bibitem{LLS} R. Lawther, M. W. Liebeck and G. M. Seitz,
\emph{Fixed point ratios in actions of finite exceptional groups of
Lie type}, Pacific J. Math. 205 (2002), 393--464.
\bibitem{BT} T. C. Burness and A. R. Thomas, \emph{On the involution
fixity of exceptional groups of Lie type}, arXiv:1707.03562v2 (2018).
\bibitem{AS} M. Aschbacher and G. M. Seitz, \emph{Involutions in
Chevalley groups over fields of even order}, Nagoya Math. J. 63
(1976), 1--91.
\bibitem{ASerr} M. Aschbacher and G. M. Seitz, \emph{Corrections to
``Involutions in Chevalley groups over fields of even order''},
Nagoya Math. J. 72 (1978), 135--136.
\end{thebibliography}
\end{document}
